\section{Introduction}

Ancient books are important carriers of cultural heritage. Converting scanned historical documents into searchable and editable text is therefore a basic requirement for preservation-oriented digitization and downstream knowledge organization. Historical document analysis has shown that writing medium, imaging quality, script style, and degradation can strongly affect recognition reliability~\cite{lombardi2020historicalsurvey}. Historical book recognition benchmarks further indicate that real collections bring persistent variation in fonts, degradation, and page appearance~\cite{antonacopoulos2013hbr}. For Chinese historical documents, vertical writing, character localization, historical glyph forms, and collection-specific character usage make recognition especially challenging~\cite{tang2020hrcenternet,wu2020precise,ma2024reading}.

Scene text recognition has developed a mature segmentation-free paradigm. Representative recognizers predict a character sequence directly from a cropped text image, using CTC, attention, language modeling, or Transformer-based sequence modeling~\cite{shi2017crnn,shi2016rare,baek2019wrong,fang2021abinet,du2022svtr,bautista2022parseq}. Vertical ancient-text recognition shares the same input-output form: given a cropped line or single-column image, the model predicts the corresponding character sequence. The visual and lexical conditions, however, differ substantially. Paper stains, ink diffusion, low contrast, broken strokes, and missing local evidence weaken character appearance; historical glyph variants and collection-specific rare characters enlarge the output space. This paper focuses on cropped vertical single-column recognition. This setting requires line-level images and transcriptions, and is more compatible with large-scale digitization than workflows that require dense character-level segmentation at inference time.

Segmentation-free sequence recognition is attractive because training can rely on line-level transcriptions without explicit character boxes or alignments. Yet characters in vertical ancient text are not only sequence tokens; they are also visual instances arranged along a vertical spatial axis. Their relative positions, scales, and local spacing carry structural information. Line-level CTC supervision constrains the final transcription, but it does not directly specify where each character appears in the column. When image degradation, spacing variation, and vertical organization interact, purely implicit sequence learning may underuse this spatial structure.

A direct way to exploit character positions is to introduce character detection. Traditional detection-recognition pipelines, however, often depend on inference-time boxes, confidence thresholds, filtering, and spatial sorting. Errors in any of these steps can propagate to the final transcription. We use character boxes differently: they act only as training-time localization supervision, so that queries learn character positions and vertical organization, while final recognition remains line-level CTC sequence prediction. This allows the model to benefit from character-level spatial signals without changing the inference interface.

Another practical issue is character-set mismatch. Ancient-text datasets often share common characters but include different rare characters, variants, or collection-specific symbols. If a target dataset uses a different charset, randomly initializing the classifier discards learned classifier parameters for shared characters, while reusing the old classifier cannot cover newly introduced characters. Character-set change is therefore not only an output-dimension issue; it also affects how visual-character correspondences are inherited.

We propose SAQT, a structure-aware query learning method for vertical ancient Chinese text recognition. SAQT uses detection-style queries to learn character-localization information and converts the sorted query outputs into a CTC sequence. When the target charset differs from the reference charset associated with the localization-aware query representation, charset-aware classifier adaptation inherits shared-character classifier rows and initializes rows for new characters. At inference time, the same recognizer can be evaluated with direct greedy decoding or with a lightweight blank/nonblank bias fixed on the development split. The paper makes the following contributions:

\begin{itemize}
  \item We introduce a structure-aware query learning method that uses character boxes as intermediate training supervision rather than final inference outputs, enabling query representations to capture character-localization information and vertical spatial organization.
  \item We give a charset-aware classifier adaptation strategy that inherits classifier parameters for shared characters and initializes target-only characters from source classifier rows, providing a training initialization mechanism for adapting the query model to different output spaces.
  \item We evaluate SAQT on MTHv2 and HDRC with a unified AR/CR protocol, and compare it with multiple vertical-input-adapted scene text recognition baselines considered in this paper.
\end{itemize}
